\documentclass[11pt,a4paper]{article}

\usepackage[utf8]{inputenc}
\usepackage[T1]{fontenc}
\usepackage[polish]{babel}
\usepackage{lmodern}
\usepackage[a4paper,margin=2.2cm,headheight=14pt]{geometry}
\usepackage{booktabs}
\usepackage{tabularx}
\usepackage{array}
\usepackage[table]{xcolor}
\usepackage{graphicx}
\usepackage{hyperref}
\usepackage{fancyhdr}
\usepackage{enumitem}
\usepackage{float}

\hypersetup{
  colorlinks=true,
  linkcolor=blue!45!black,
  urlcolor=blue!45!black,
  pdftitle={Podstawy TCP/IP i skanowanie sieci za pomocą Nmap},
  pdfauthor={},
  pdfsubject={Sprawozdanie z ćwiczenia laboratoryjnego}
}

\definecolor{navy}{HTML}{183153}
\definecolor{pale}{HTML}{EAF0F6}
\definecolor{barblue}{HTML}{3973AC}
\definecolor{muted}{HTML}{555555}

\pagestyle{fancy}
\fancyhf{}
\fancyhead[L]{\small Sieci komputerowe -- ćwiczenie 4}
\fancyhead[R]{\small Nmap}
\fancyfoot[C]{\thepage}
\setlength{\parskip}{0.45em}
\setlength{\parindent}{0pt}
\renewcommand{\arraystretch}{1.2}
\setlist[itemize]{leftmargin=1.5em,itemsep=0.15em,topsep=0.2em}
\setlist[enumerate]{leftmargin=1.7em,itemsep=0.2em,topsep=0.2em}

\newcommand{\statbar}[3]{%
  \texttt{#1} & \rule{#2cm}{0.85ex} & #3\\
}
\newcommand{\osbar}[3]{%
  #1 & \rule{#2cm}{0.85ex} & #3\\
}

\begin{document}

\begin{titlepage}
\thispagestyle{empty}
\vspace*{0.8cm}
\begin{tabularx}{\textwidth}{|>{\raggedright\arraybackslash}p{0.19\textwidth}|X|X|}
\hline
\textit{Grupa lab.} & \textit{Data wykonania} & \textit{Data odbioru}\\[1.2cm]
\hline
\multicolumn{3}{|l|}{\textit{Temat ćwiczenia: Podstawy TCP/IP, adresowania IP i skanowania sieci za pomocą Nmap}}\\[1.2cm]
\hline
\multicolumn{2}{|p{0.62\textwidth}|}{\textit{Imiona i nazwiska:} \vspace{0.7cm}} &
\textit{Ocena i uwagi} \vspace{0.7cm}\\
\hline
\end{tabularx}

\vspace{1.6cm}
\begin{center}
{\Large\bfseries\color{navy} SPRAWOZDANIE Z ĆWICZENIA LABORATORYJNEGO}\\[0.5cm]
{\huge\bfseries Podstawy TCP/IP i skanowanie sieci za pomocą Nmap}\\[0.6cm]
{\large Część praktyczna}
\end{center}

\vfill
\begin{tabularx}{\textwidth}{@{}lX@{}}
\textbf{Data wykonania pomiarów:} & 25 listopada 2026 r.\\
\textbf{System i narzędzie:} & Ubuntu/Linux, Nmap 7.94SVN\\
\textbf{Zakres opracowania:} & rozpoznawanie hostów, mapowanie DNS, skanowanie TCP SYN, rozpoznawanie usług i systemów operacyjnych\\
\end{tabularx}
\vspace{1.0cm}
\end{titlepage}

\section{Cel i zakres ćwiczenia}
Celem ćwiczenia było praktyczne poznanie sposobów rozpoznawania hostów i usług w sieci
z wykorzystaniem programu Nmap oraz powiązanie obserwowanych wyników z podstawami
adresowania IPv4, protokołami TCP/IP i stanami portów. Wykonaliśmy skanowania listujące
adresy i nazwy DNS (\texttt{-sL}), wykrywanie aktywnych hostów (\texttt{-sn}), skanowanie
TCP SYN (\texttt{-sS}) oraz wybrane skanowania usług i systemów operacyjnych
(\texttt{-sV}, \texttt{-O} i \texttt{-A}).

Opracowanie obejmuje dane z dwóch zestawów pomiarowych: sieci prywatnej
\texttt{192.168.102.0/24} oraz sześciu podsieci z zakresu \texttt{149.156.111.0/24}--
\texttt{149.156.116.0/24}. Wyniki traktujemy jako migawkę stanu sieci z chwili pomiaru.
Nie utożsamiamy nazw usług rozpoznawanych wyłącznie na podstawie numeru portu
z potwierdzoną wersją oprogramowania.

\section{Podstawy teoretyczne}
\subsection{Adresacja i porty}
Adres IPv4 identyfikuje interfejs w sieci i składa się z 32 bitów. Maska podsieci
określa, która część adresu opisuje sieć, a która hosta. W pomiarach użyto podsieci
\texttt{/24}; każda z nich obejmuje 256 adresów IPv4, z których część może być
zarezerwowana na adres sieci i rozgłoszeniowy. Port TCP jest numerowanym punktem
końcowym usługi działającej na hoście. Numer portu podpowiada możliwą usługę, ale sam
w sobie nie potwierdza jej rodzaju ani wersji.

\subsection{Zastosowane techniki}
\begin{itemize}
  \item \textbf{\texttt{-sL}} -- listowanie adresów z próbą uzyskania nazw DNS; nie
        sprawdza, czy host odpowiada, ani nie bada portów.
  \item \textbf{\texttt{-sn}} -- wykrywanie hostów bez skanowania portów. Wynik zależy
        od dostępności użytych metod wykrywania i filtrów sieciowych.
  \item \textbf{\texttt{-sS}} -- skanowanie TCP typu SYN. Odpowiedź SYN/ACK wskazuje
        zwykle port otwarty, RST -- zamknięty, a brak odpowiedzi lub komunikat
        niedostępności może skutkować klasyfikacją jako filtrowany.
  \item \textbf{\texttt{-sV}} -- dodatkowe sondowanie otwartych portów w celu
        rozpoznania usługi i, jeśli to możliwe, jej wersji.
  \item \textbf{\texttt{-O}} -- próba określenia systemu operacyjnego na podstawie
        cech odpowiedzi stosu sieciowego. Wynik jest szacunkiem, zależnym m.in. od
        liczby i stanów dostępnych portów.
  \item \textbf{\texttt{-A}} -- zestaw funkcji zwiększających szczegółowość
        rozpoznania; w zebranych wynikach obejmował m.in. wykrywanie usług i systemu.
\end{itemize}

\section{Stanowisko i przebieg}
\subsection{Sprzęt i oprogramowanie}
Ćwiczenie wykonaliśmy na komputerze pracującym pod kontrolą systemu Linux
(Ubuntu). Do pomiarów wykorzystaliśmy Nmap w wersji 7.94SVN. Wyniki zapisywaliśmy
w plikach tekstowych; część przebiegu udokumentowaliśmy zrzutami ekranu terminala.

\subsection{Przebieg prac}
\begin{enumerate}
  \item Zebraliśmy listy adresów i mapowania DNS dla wybranych podsieci za pomocą
        skanowania listującego.
  \item Wykryliśmy aktywne hosty w poszczególnych zakresach bez badania portów.
  \item Dla wybranej listy ośmiu hostów z zakresu \texttt{149.156.116.0/24}
        przeanalizowaliśmy stany portów TCP metodą SYN.
  \item Dla ośmiu wybranych hostów z sieci \texttt{192.168.102.0/24} zebraliśmy
        informacje o usługach; dla dwóch wykonaliśmy również szczegółowe skanowanie
        typu \texttt{-A}.
  \item Dla ośmiu wybranych hostów z zakresu \texttt{149.156.111.0/24} wykonaliśmy
        próby rozpoznania systemu operacyjnego.
  \item Zestawiliśmy liczbę aktywnych hostów, zaobserwowane porty i usługi,
        a następnie oceniliśmy ograniczenia wnioskowania.
\end{enumerate}

\subsection{Schemat ćwiczenia}
\begin{center}
\small
\begin{tabular}{c@{\hspace{0.15em}$\longrightarrow$\hspace{0.15em}}c@{\hspace{0.15em}$\longrightarrow$\hspace{0.15em}}c@{\hspace{0.15em}$\longrightarrow$\hspace{0.15em}}c}
\fbox{\parbox[c][1.15cm][c]{0.18\textwidth}{\centering 1. Listowanie\\\texttt{-sL}}} &
\fbox{\parbox[c][1.15cm][c]{0.18\textwidth}{\centering 2. Wykrywanie hostów\\\texttt{-sn}}} &
\fbox{\parbox[c][1.15cm][c]{0.18\textwidth}{\centering 3. Stany portów TCP\\\texttt{-sS}}} &
\fbox{\parbox[c][1.15cm][c]{0.18\textwidth}{\centering 4. Usługi i system\\\texttt{-sV}, \texttt{-O}}}
\end{tabular}
\end{center}

\begin{table}[H]
\centering
\caption{Zakresy i rodzaje dostępnych pomiarów}
\begin{tabularx}{\textwidth}{>{\ttfamily}l X X}
\toprule
\normalfont\textbf{Zakres} & \textbf{Zebrane dane} & \textbf{Uwagi}\\
\midrule
192.168.102.0/24 & \texttt{-sL}, \texttt{-sn}; wybrane \texttt{-sV} i \texttt{-A} &
61 hostów wykrytych jako aktywne; usługi zbadano na wybranych hostach.\\
149.156.111.0/24 & \texttt{-sL}, \texttt{-sn}; wybrane \texttt{-O} &
26 hostów aktywnych; osiem raportów prób identyfikacji systemu.\\
149.156.112.0/24 & \texttt{-sL}, \texttt{-sn} &
20 hostów aktywnych.\\
149.156.113.0/24 & \texttt{-sL}, \texttt{-sn} &
14 hostów aktywnych.\\
149.156.114.0/24 & \texttt{-sL}, \texttt{-sn} &
37 hostów aktywnych.\\
149.156.115.0/24 & \texttt{-sL}, \texttt{-sn} &
22 hosty aktywne.\\
149.156.116.0/24 & \texttt{-sL}, \texttt{-sn}; wybrane \texttt{-sS} &
13 hostów aktywnych; skan SYN dotyczył ośmiu wybranych hostów.\\
\bottomrule
\end{tabularx}
\end{table}

\section{Wyniki pomiarów}
\subsection{Wykrywanie hostów i mapowanie DNS}
W sześciu analizowanych podsieciach publicznych wykryliśmy łącznie 132 aktywne
adresy spośród 1536 sprawdzonych adresów (około
8{,}6\%). W sieci prywatnej \texttt{192.168.102.0/24} wykryliśmy 61 hostów spośród
256 adresów (około 23{,}8\%). Wartości te opisują wyłącznie odpowiedzi rozpoznane
przez Nmap w chwili skanowania, a nie pełną liczbę urządzeń fizycznie obecnych
w sieci.

Wyniki uzyskaliśmy poleceniami wykrywania hostów dla kolejnych podsieci
\texttt{/24}; odpowiadają im pliki \texttt{c\_sn.txt} oraz
\texttt{c\_112sn.txt}--\texttt{c\_116sn.txt}. W każdym skanowaniu sprawdzono 256
adresów. Dla przykładu w \texttt{149.156.116.0/24} Nmap zgłosił 13 hostów
aktywnych w czasie 4{,}34 s, a w \texttt{192.168.102.0/24} -- 61 hostów
w czasie 1{,}59 s. Polecenia listujące DNS odpowiadają plikom
\texttt{c\_sL.txt} oraz \texttt{c\_112sL.txt}--\texttt{c\_116sL.txt}.
W trybie \texttt{-sL} Nmap wypisał mapowania dla badanych adresów, ale
podsumowanie wynosiło 0 aktywnych hostów, ponieważ ten tryb nie wykonuje
sprawdzenia dostępności.

\begin{table}[H]
\centering
\caption{Przykładowe mapowania DNS odczytane z listowania podsieci
\texttt{149.156.112.0/24}}
\begin{tabular}{ll}
\toprule
\textbf{Adres IP} & \textbf{Nazwa zwrócona przez DNS}\\
\midrule
\texttt{149.156.112.6} & \texttt{krywan.metal.agh.edu.pl}\\
\texttt{149.156.112.13} & \texttt{umbreon.metal.agh.edu.pl}\\
\texttt{149.156.112.15} & \texttt{metal.metal.agh.edu.pl}\\
\texttt{149.156.112.16} & \texttt{cmms-edit.agh.edu.pl}\\
\texttt{149.156.112.42} & \texttt{willow.metal.agh.edu.pl}\\
\texttt{149.156.112.124} & \texttt{www.futurestudent.agh.edu.pl}\\
\bottomrule
\end{tabular}
\end{table}

\begin{table}[H]
\centering
\caption{Liczba aktywnych hostów wykrytych w skanowaniu \texttt{-sn}}
\begin{tabular}{lrr}
\toprule
\textbf{Podsieć} & \textbf{Aktywne hosty} & \textbf{Udział w 256 adresach}\\
\midrule
\texttt{149.156.111.0/24} & 26 & 10{,}2\%\\
\texttt{149.156.112.0/24} & 20 & 7{,}8\%\\
\texttt{149.156.113.0/24} & 14 & 5{,}5\%\\
\texttt{149.156.114.0/24} & 37 & 14{,}5\%\\
\texttt{149.156.115.0/24} & 22 & 8{,}6\%\\
\texttt{149.156.116.0/24} & 13 & 5{,}1\%\\
\midrule
\textbf{Suma sześciu podsieci} & \textbf{132} & \textbf{8{,}6\% z 1536}\\
\texttt{192.168.102.0/24} & 61 & 23{,}8\%\\
\bottomrule
\end{tabular}
\end{table}

Wyniki \texttt{-sL} zawierają odpowiedzi resolvera dla części adresów i nazwy
kanoniczne hostów. Podsumowanie Nmapa dla tego trybu podaje zero aktywnych hostów;
jest to oczekiwane, ponieważ \texttt{-sL} nie wykonuje testu dostępności. Brak
nazwy DNS przy adresie nie oznacza, że adres jest nieaktywny.

\subsection{Skanowanie TCP SYN ośmiu hostów}
Polecenie widoczne na zrzucie z terminala miało postać
\texttt{sudo nmap -sS -iL lista\_skan\_sS.txt}. Dla ośmiu hostów z listy
pomiarowej w podsieci \texttt{149.156.116.0/24} raporty
wykazały 22 otwarte porty TCP i trzy jawnie wypisane porty filtrowane. Dodatkowe
podsumowania Nmapa wskazują filtrowanie znacznej części skanowanych portów na
czterech z ośmiu hostów. Na pozostałych hostach przeważały porty zamknięte.
W tabeli częstości poniżej uwzględniliśmy wszystkie 22 otwarte porty, nie zaś
liczbę unikatowych nazw hostów.

Pełny wynik zapisano w pliku \texttt{skan\_listy\_8el.txt}. Tabela przedstawia
stany portów host po hoście, dzięki czemu widoczne jest, które usługi
odpowiadały na sondy, a które były filtrowane.

\begin{table}[H]
\centering
\small
\caption{Wyniki skanowania SYN zapisane w \texttt{skan\_listy\_8el.txt}}
\begin{tabularx}{\textwidth}{>{\ttfamily}l X}
\toprule
\normalfont\textbf{Host} & \textbf{Wynik: otwarte porty / inne stany}\\
\midrule
149.156.116.12 &
21/tcp \textit{open} (ftp), 80/tcp \textit{open} (http); pozostałe 998 portów zamknięte.\\
149.156.116.17 &
21/tcp \textit{open} (ftp), 80/tcp \textit{open} (http); pozostałe 998 portów zamknięte.\\
149.156.116.76 &
21/tcp \textit{open} (ftp); pozostałe 999 portów filtrowane (bez odpowiedzi).\\
149.156.116.129 &
135/tcp (msrpc), 139/tcp (netbios-ssn) i 445/tcp (microsoft-ds) \textit{filtered}; pozostałe 997 portów zamknięte.\\
149.156.116.133 &
22/tcp (ssh), 23/tcp (telnet), 80/tcp (http) \textit{open}; pozostałe 997 portów filtrowane.\\
149.156.116.136 &
21/tcp (ftp), 23/tcp (telnet), 80/tcp (http), 139/tcp (netbios-ssn), 514/tcp (shell), 515/tcp (printer), 631/tcp (ipp), 7443/tcp (oracleas-https), 9100/tcp (jetdirect) \textit{open}; pozostałe 991 zamknięte.\\
149.156.116.137 &
53/tcp (domain), 1723/tcp (pptp), 2000/tcp (cisco-sccp), 8291/tcp (unknown) \textit{open}; pozostałe 996 zamknięte.\\
149.156.116.199 &
21/tcp \textit{open} (ftp); pozostałe 999 portów filtrowane (bez odpowiedzi).\\
\bottomrule
\end{tabularx}
\end{table}

\begin{table}[H]
\centering
\caption{Najczęściej rozpoznane usługi w ośmiu raportach \texttt{-sS}}
\begin{tabular}{lrr}
\toprule
\textbf{Usługa (etykieta Nmapa)} & \textbf{Liczba otwartych portów} & \textbf{Udział w 22}\\
\midrule
\texttt{ftp} & 5 & 22{,}7\%\\
\texttt{http} & 4 & 18{,}2\%\\
\texttt{telnet} & 2 & 9{,}1\%\\
\texttt{ssh} & 1 & 4{,}5\%\\
\texttt{netbios-ssn} & 1 & 4{,}5\%\\
\texttt{shell} & 1 & 4{,}5\%\\
\texttt{printer} & 1 & 4{,}5\%\\
\texttt{ipp} & 1 & 4{,}5\%\\
\texttt{oracleas-https} & 1 & 4{,}5\%\\
\texttt{jetdirect} & 1 & 4{,}5\%\\
\texttt{domain} & 1 & 4{,}5\%\\
\texttt{pptp} & 1 & 4{,}5\%\\
\texttt{cisco-sccp} & 1 & 4{,}5\%\\
\texttt{unknown} & 1 & 4{,}5\%\\
\bottomrule
\end{tabular}
\end{table}

\begin{figure}[H]
\centering
\begin{tabular}{l l r}
\textbf{Usługa} & \textbf{Liczba otwartych portów} & \textbf{n}\\
\statbar{ftp}{2.5}{5}
\statbar{http}{2.0}{4}
\statbar{telnet}{1.0}{2}
\statbar{pozostałe}{0.5}{1 każda}
\end{tabular}
\caption{Wykres częstości usług w wybranym zbiorze ośmiu hostów. Każda z pozostałych
dziesięciu etykiet wystąpiła jeden raz.}
\end{figure}

\subsection{Wybrane skanowania usług i systemów operacyjnych}
W ośmiu skanowaniach usług hostów z sieci prywatnej odnotowaliśmy łącznie 39
otwartych portów. Najczęściej Nmap wskazywał \texttt{msrpc} (12 razy) i
\texttt{flexlm} (6 razy), następnie \texttt{ssh} oraz \texttt{netbios-ssn}
(po 4 razy), a także \texttt{microsoft-ds} i \texttt{http} (po 3 razy).
Są to częstości etykiet usług w wynikach, a nie liczba unikatowych hostów.
Poniżej zestawiliśmy konkretny wynik każdego skanowania z plików
\texttt{sV150.txt}, \texttt{sV165.txt}, \texttt{sV196.txt},
\texttt{sV199.txt}, \texttt{sV247.txt}, \texttt{sV249.txt},
\texttt{sV250.txt} i \texttt{sV253.txt}. Wersje oprogramowania podajemy
tylko tam, gdzie Nmap faktycznie je rozpoznał.

\begin{table}[H]
\centering
\small
\caption{Otwarte porty i rozpoznane usługi w ośmiu skanowaniach
\texttt{-sV} w sieci prywatnej}
\begin{tabularx}{\textwidth}{>{\ttfamily}l X}
\toprule
\normalfont\textbf{Host} & \textbf{Otwarte porty TCP i identyfikacja}\\
\midrule
192.168.102.150 &
22/ssh -- OpenSSH 9.6p1, Ubuntu; system Linux.\\
192.168.102.165 &
135/msrpc, 139/netbios-ssn, 445/microsoft-ds, 1947/http (HASP LM 17.00), 3389/ms-wbt-server, 25734/flexlm, 25735/flexlm, 27000/flexlm.\\
192.168.102.196 &
22/ssh -- OpenSSH 9.6p1, Ubuntu; system Linux.\\
192.168.102.199 &
22/ssh -- OpenSSH 9.6p1, Ubuntu; system Linux.\\
192.168.102.247 &
135/msrpc, 139/netbios-ssn, 445/microsoft-ds (niepewna etykieta), 5357/http -- Microsoft HTTPAPI 2.0.\\
192.168.102.249 &
135/msrpc, 139/netbios-ssn, 445/microsoft-ds, 1025/msrpc, 1026/msrpc, 27000/flexlm.\\
192.168.102.250 &
135/msrpc, 139/netbios-ssn, 445/microsoft-ds, 3389/ms-wbt-server, 5666/tcpwrapped, 9102/jetdirect, 27000/flexlm, 49152/msrpc, 49153/msrpc, 49156/msrpc, 49163/msrpc, 49165/msrpc, 49167/msrpc, 49176/flexlm.\\
192.168.102.253 &
22/ssh -- OpenSSH 7.4; 80/http -- Apache httpd 2.4.6 (CentOS); 111/rpcbind; 6000/X11 (dostęp zablokowany).\\
\bottomrule
\end{tabularx}
\end{table}

W dwóch szczegółowych raportach z sieci prywatnej Nmap rozpoznał systemy
z rodziny Microsoft Windows Server: dla adresu \texttt{192.168.102.165}
Windows Server 2019, a dla \texttt{192.168.102.250} Windows Server 2008 R2.
W drugim przypadku część ogólnego dopasowania systemu była niejednoznaczna;
szczegółowa informacja pojawiła się w wynikach rozpoznawania SMB. Z tego powodu
traktujemy te identyfikacje jako wskazówki zebrane w konkretnym pomiarze, a nie
niezależnie potwierdzoną inwentaryzację systemów.

\begin{table}[H]
\centering
\small
\caption{Dodatkowe ustalenia z dwóch szczegółowych skanowań \texttt{-A}}
\begin{tabularx}{\textwidth}{>{\ttfamily}l X}
\toprule
\normalfont\textbf{Host} & \textbf{Wynik skanowania i szczegóły rozpoznania}\\
\midrule
192.168.102.165 &
8 otwartych portów, 992 zamknięte; Windows Server 2019 Standard 17763, nazwa komputera CHARIZARD. Rozpoznano m.in. serwer licencji HASP LM 17.00 na 1947/tcp, RDP na 3389/tcp oraz FlexLM na 25734/tcp, 25735/tcp i 27000/tcp. Czas skanowania: 31,14 s.\\
192.168.102.250 &
14 otwartych portów, 986 zamkniętych; szczegółowa identyfikacja SMB wskazała Windows Server 2008 R2 Enterprise 7600, nazwę Procast-SIM i grupę WORKGROUP. Wykryto m.in. RPC, RDP, FlexLM oraz porty dynamiczne RPC. Czas skanowania: 93,16 s.\\
\bottomrule
\end{tabularx}
\end{table}

\begin{figure}[H]
\centering
\begin{tabular}{l l r}
\textbf{Usługa} & \textbf{Liczba rozpoznań} & \textbf{n}\\
\statbar{msrpc}{2.4}{12}
\statbar{flexlm}{1.2}{6}
\statbar{ssh}{0.8}{4}
\statbar{netbios-ssn}{0.8}{4}
\statbar{microsoft-ds}{0.6}{3}
\statbar{http}{0.6}{3}
\end{tabular}
\caption{Najczęstsze etykiety usług w ośmiu skanowaniach \texttt{-sV}
w sieci \texttt{192.168.102.0/24}. Pozostałe etykiety wystąpiły łącznie
siedem razy.}
\end{figure}

\begin{table}[H]
\centering
\caption{Wynik identyfikacji systemu w dwóch szczegółowych skanowaniach}
\begin{tabular}{lll}
\toprule
\textbf{Adres} & \textbf{Rozpoznany system} & \textbf{Ocena wyniku}\\
\midrule
\texttt{192.168.102.165} & Windows Server 2019 & szczegółowe dopasowanie Nmap/SMB\\
\texttt{192.168.102.250} & Windows Server 2008 R2 & część dopasowania ogólnego niejednoznaczna\\
\bottomrule
\end{tabular}
\end{table}

\begin{figure}[H]
\centering
\begin{tabular}{l l r}
\textbf{Rozpoznany system} & \textbf{Liczba hostów} & \textbf{n}\\
\osbar{Windows Server 2019}{1.0}{1}
\osbar{Windows Server 2008 R2}{1.0}{1}
\end{tabular}
\caption{Systemy rozpoznane w dwóch szczegółowych skanowaniach w sieci prywatnej.
Próba jest zbyt mała, aby wnioskować o udziale systemów w całej sieci.}
\end{figure}

Dla ośmiu hostów z podsieci \texttt{149.156.111.0/24} zachowały się raporty
\texttt{-O} w plikach \texttt{c\_wyn1.txt}--\texttt{c\_wyn8.txt}. W większości
przypadków Nmap nie wskazał jednego pewnego systemu: w raportach powtarza się
przypuszczenie o urządzeniu typu Digi PortServer (85--90\%) oraz alternatywne
dopasowania Windows lub FreeBSD. Tylko ostatni raport nie zawierał wyniku
identyfikacji OS. Szczegóły widoczne w raportach zebraliśmy poniżej.

\begin{table}[H]
\centering
\scriptsize
\caption{Porty i przypuszczenia dotyczące systemu w wynikach \texttt{-O}
dla ośmiu hostów publicznych}
\begin{tabularx}{\textwidth}{>{\ttfamily}l X X}
\toprule
\normalfont\textbf{Host} & \textbf{Otwarte porty TCP} & \textbf{Wynik OS / filtrowanie}\\
\midrule
149.156.111.104 &
21/ftp, 443/https &
Windows 7 / Server 2008 (87\%); bez dokładnego dopasowania; 998 filtrowanych.\\
149.156.111.149 &
21/ftp, 22/ssh; 8000/tcp zamknięty &
FreeBSD 6.2 (88\%) lub Windows Server 2008 (87\%); bez dokładnego dopasowania; 997 filtrowanych.\\
149.156.111.212 &
21/ftp &
Digi PortServer (90\%); bez dokładnego dopasowania; 999 filtrowanych.\\
149.156.111.53 &
21/ftp, 22/ssh, 80/http, 113/ident, 443/https &
FreeBSD 6.2 (88\%); bez dokładnego dopasowania; 991 filtrowanych.\\
149.156.111.36 &
21/ftp, 80/http, 443/https &
Digi PortServer (90\%); bez dokładnego dopasowania; 997 filtrowanych.\\
149.156.111.191 &
21/ftp &
Digi PortServer (90\%); bez dokładnego dopasowania; 999 filtrowanych.\\
149.156.111.172 &
21/ftp &
Digi PortServer (90\%); bez dokładnego dopasowania; 999 filtrowanych.\\
149.156.111.162 &
21/ftp, 80/http, 99/metagram, 5000/upnp &
Brak wyniku identyfikacji OS; 996 portów zamkniętych.\\
\bottomrule
\end{tabularx}
\end{table}

Podane procenty oznaczają poziom dopasowania hipotezy Nmapa, a nie prawdopodobieństwo
potwierdzone niezależnym pomiarem. W siedmiu raportach systemu nie udało się
zidentyfikować jednoznacznie, m.in. z powodu licznych portów filtrowanych.

\section{Opracowanie i interpretacja wyników}
\begin{enumerate}
  \item Najwięcej hostów spośród sześciu analizowanych podsieci wykryliśmy w sieci
        \texttt{149.156.114.0/24} (37), najmniej w \texttt{149.156.116.0/24}
        (13). Różnice pokazują, że liczba aktywnych adresów istotnie zależy od
        konkretnego segmentu.
  \item Skanowanie listujące i wykrywanie hostów odpowiadają na różne pytania.
        \texttt{-sL} pozwala zestawiać adresy z nazwami, ale nie potwierdza
        dostępności. \texttt{-sn} szacuje liczbę odpowiadających hostów, ale
        odpowiedź może być blokowana przez zaporę albo zależna od bieżącego stanu
        urządzenia.
  \item W próbce ośmiu hostów poddanych \texttt{-sS} najczęstszą etykietą usługi
        było \texttt{ftp}, a następnie \texttt{http}. Zarejestrowano również
        \texttt{telnet} oraz usługi związane m.in. z drukowaniem i zarządzaniem
        urządzeniami. Ze względu na małą, celowo wybraną próbę nie ekstrapolujemy
        tych proporcji na całą podsieć.
  \item Na czterech z ośmiu badanych hostów Nmap raportował filtrowanie części
        portów. Stan filtrowany oznacza, że skaner nie uzyskał odpowiedzi
        pozwalającej rozstrzygnąć, czy port jest otwarty, czy zamknięty; nie jest
        równoznaczny z potwierdzoną luką ani awarią.
  \item Rozpoznawanie usług przez \texttt{-sV} dostarczyło więcej informacji
        niż samo przypisanie nazwy do numeru portu. Identyfikacja systemu
        operacyjnego była mniej pewna i zależała od warunków skanowania oraz
        dostępnych odpowiedzi hosta.
\end{enumerate}

\subsection{Ograniczenia i kompletność materiału}
W przekazanym zbiorze nie ma pliku z wynikami UDP, przechwycenia \texttt{p0f}
ani parametrów konfiguracji interfejsu komputera (adresu, maski, bramy i DNS).
Nie uzupełniamy tych pozycji na podstawie przypuszczeń. Skanowania usług i systemów
wykonaliśmy tylko dla wybranych hostów; zatem nie są one pełną inwentaryzacją
żadnej z podsieci. Zestawy \texttt{-sV}, \texttt{-A}, \texttt{-O} i \texttt{-sS}
różnią się zakresem oraz celem, dlatego ich wyników nie sumujemy w jedną wspólną
statystykę.

\section{Wnioski}
\begin{enumerate}
  \item Zrealizowane pomiary pozwoliły nam przećwiczyć rozdzielenie trzech zadań:
        listowania adresów i DNS, wykrywania aktywności hostów oraz badania
        portów i usług.
  \item W sześciu zebranych podsieciach \texttt{/24} wykryliśmy łącznie 132
        aktywne adresy, natomiast w prywatnej sieci \texttt{192.168.102.0/24}
        -- 61. Są to liczby hostów odpowiadających na zastosowane sondy, a nie
        bezwarunkowo pełne liczebności sieci.
  \item W ośmiu raportach SYN zarejestrowaliśmy 22 otwarte porty; najczęściej
        występowały etykiety \texttt{ftp} i \texttt{http}. Jednocześnie cztery
        hosty wykazywały filtrowanie, co potwierdza, że interpretując brak
        odpowiedzi, należy odróżniać stan \textit{filtered} od
        \textit{closed}.
  \item Skanowanie wersji usług i identyfikacja systemu operacyjnego uzupełniają
        obraz hosta, ale wyniki Nmapa pozostają estymacją. W szczególności
        niejednoznacznych dopasowań OS nie należy przedstawiać jako pewnych.
  \item Wnioski o popularności usług odnoszą się do małej, wybranej próbki.
        Nie przenosimy ich na całą podsieć. W następnym uzupełnieniu raportu
        dołączymy skany podpisanych notatek z zajęć oraz, jeśli będą dostępne,
        brakujące wyniki i dane konfiguracyjne.
\end{enumerate}

\clearpage
\section*{Załącznik A. Zrzuty ekranu z ćwiczenia}
Poniższe zrzuty dokumentują kolejne etapy pracy z Nmapem. Zestawiliśmy je
z odpowiednimi plikami tekstowymi: obrazy pokazują polecenia i fragmenty
interaktywnego wyniku, natomiast kompletne liczby hostów i portów w tabelach
przepisaliśmy z zapisanych raportów.

\begin{figure}[H]
\centering
\includegraphics[width=\textwidth,height=0.36\textheight,keepaspectratio]
  {\detokenize{dane_bedkowski/zdjecia_ss/Zrzut ekranu z 2026-11-25 12-43-02.png}}
\caption{Skanowanie podsieci \texttt{149.156.115.0/24}. U góry widzimy
podsumowanie wcześniejszego wykrywania hostów: 22 adresy aktywne z 256.
Niżej uruchomiono listowanie \texttt{-sL} i zapis wyniku do
\texttt{c\_115sL.txt}; wypisane nazwy, np. \texttt{emi.wibro.agh.edu.pl},
pochodzą z mapowania DNS.}
\end{figure}
\begin{figure}[H]
\centering
\includegraphics[width=\textwidth,height=0.36\textheight,keepaspectratio]
  {\detokenize{dane_bedkowski/zdjecia_ss/Zrzut ekranu z 2026-11-25 12-43-19.png}}
\caption{Dwa uzupełniające się wyniki dla \texttt{149.156.116.0/24}.
Listowanie DNS kończy się komunikatem \texttt{0 hosts up}, co nie oznacza,
że adresy są nieaktywne. Następne polecenie \texttt{-sn} wykrywa 13
odpowiadających hostów; widoczne nazwy to m.in. \texttt{dzk.ceramika.agh.edu.pl}
i \texttt{a3-09.gorn.agh.edu.pl}.}
\end{figure}
\begin{figure}[H]
\centering
\includegraphics[width=\textwidth,height=0.36\textheight,keepaspectratio]
  {\detokenize{dane_bedkowski/zdjecia_ss/Zrzut ekranu z 2026-11-25 12-43-30.png}}
\caption{Powtórne wyświetlenie wykrywania aktywnych hostów i listowania DNS
w podsieci \texttt{149.156.116.0/24}. Na dole widoczna jest nieudana próba
użycia opcji \texttt{-iL} z adresem IP zamiast ścieżki do pliku listy.
Komunikat o błędzie otwarcia pliku oznacza, że ta próba nie jest wynikiem
skanowania hosta i nie została włączona do opracowania.}
\end{figure}
\begin{figure}[H]
\centering
\includegraphics[width=\textwidth,height=0.36\textheight,keepaspectratio]
  {\detokenize{dane_bedkowski/zdjecia_ss/Zrzut ekranu z 2026-11-25 12-43-39.png}}
\caption{Skan SYN listy zapisanej w \texttt{lista\_skan\_sS.txt}. Najpierw
pokazano wynik dla \texttt{149.156.116.12}: otwarte porty 21/tcp (FTP) i
80/tcp (HTTP). Następnie widoczne jest uruchomienie skanowania listy;
na kolejnych hostach obserwujemy m.in. porty filtrowane na
\texttt{149.156.116.129}. Pełne wyniki wszystkich ośmiu hostów znajdują się
w tabeli opracowanej z \texttt{skan\_listy\_8el.txt}.}
\end{figure}

\clearpage
\section*{Załącznik B. Podpisane notatki z zajęć}
\begin{center}
\vspace*{3cm}
\fcolorbox{navy}{pale}{%
  \parbox{0.82\textwidth}{%
    \centering
    \vspace{1.2cm}
    \textbf{\Large Miejsce na skany podpisanych notatek}\\[0.6cm]
    Skany zostaną dołączone na ostatniej stronie po ich otrzymaniu.\\
    Na tym etapie załącznik jest celowo pozostawiony jako miejsce do uzupełnienia.
    \vspace{1.2cm}
  }%
}
\end{center}

\end{document}
